\documentclass[12pt]{article}
\usepackage[utf8]{inputenc}
\usepackage{ctex}
\usepackage{amsmath}
\usepackage{amssymb}

\title{定义 $R$, $H$, Position}
\author{}
\date{}

\begin{document}

\maketitle

\section*{矩形区域 $R$ 的定义}

设 $\mathbf{R}^n$ 中的矩形区域 $R$ 定义为
\[
R = [a, b] = \{x \in \mathbf{R}^n \mid a \leq x \leq b\}
\]
其中 $a, b \in \mathbf{R}^n$ 且 $a_i < b_i$ ($i = 1, 2, \dots, n$)。

\section*{盒子 $H$ 的定义}

设 $H$ 为 $\mathbf{R}^{n+m}$ 中的一个"盒子"区域，定义为
\[
H = [L, L + H_E] = \{(x, y) \mid x \in \mathbf{R}^n, y \in \mathbf{R}^m, \; x \leq y \leq x + \delta\}
\]
其中 $L \in \mathbf{R}^n$ 为起始位置，$H_E \in \mathbf{R}^n$ 为边长向量（$H_E > 0$），$x$ 和 $y$ 满足 $x, y \in \mathbf{R}^n$。

\section*{位置函数 $\mathrm{Position}(x, y, P)$}

设 $x, y \in \mathbf{R}^n$ 为两个点，$P \subset \mathbf{R}^n$ 为一个超矩形区域。定义
\[
\mathrm{Position}(x, y, P) = \begin{cases}
\text{内部}, & \text{若 } x \in P \text{ 且 } y \in P \\[6pt]
\text{边界}, & \text{若 } x \in P \text{ 且 } y \notin P \\[6pt]
\text{外部}, & \text{若 } x \notin P
\end{cases}
\]

其中 $x \in P$ 表示点 $x$ 落在区域 $P$ 内部或边界上。

\end{document}